\section{Background and related work}\label{sec:background}

\subsection{The memistor (1960)}
Widrow and Hoff introduced the ADALINE and the least-mean-squares (LMS) rule in 1960~\cite{widrow1960switching}. To store its
weights they built the memistor~\cite{widrow1960adaline}. A pencil-lead (graphite) substrate sits in a copper-sulphate plating
bath. Copper is plated onto it or removed from it by a current through a third electrode, and the conductance between the
lead's two ends is sensed with a small AC signal that does not disturb the state. Widrow, Pierce and Angell described elements
with ``thousands'' of analog levels that spanned roughly 100 to 1 ohm in about ten seconds at a few milliamperes; this description
is reported in secondary sources, see~\cite{mouttet2010}. Mouttet analyses the memistor as a three-terminal memristive system
and summarises Widrow's ideal: a conductance linear in the plating charge, reversible, non-volatile and free of
hysteresis~\cite{mouttet2010}. Widrow and Lehr recount that the Madaline~I machine was built in hardware with memistor weights and used for pattern
recognition research~\cite{widrow1990}.
Write noise, up/down asymmetry, retention and endurance were never reported in a form we could find, so in our catalog these
are estimates (Section~\ref{sec:catalog}). The memistor predates Chua's memristor~\cite{chua1971} and its later physical
identification in titanium-oxide films~\cite{strukov2008}. It is the conceptual ancestor of today's three-terminal synaptic
transistors.

\subsection{Device requirements for training}
Gokmen and Vlasov~\cite{gokmen2016} defined the resistive processing unit (RPU). Updates are stochastic pulse trains whose
coincidences change the conductance by a fixed step (stochastic computing,~\cite{gaines1967}). They derived device
specifications for plain SGD on MNIST at a small error penalty: about a thousand levels, up/down asymmetry of a few percent, and
tolerance to large cycle-to-cycle and device-to-device variation. Agarwal et al.~\cite{agarwal2016} reached similar conclusions;
the requirements were extended to convolutional networks~\cite{gokmen2017cnn}, and reviews cover the device
landscape~\cite{burr2017,ielmini2018,sebastian2020,haensch2019}. Mixed-precision schemes~\cite{nandakumar2020} accumulate the
gradient digitally and program the devices rarely. They relax the device demands but give up the fully parallel in-memory update.

\subsection{Algorithms that tolerate asymmetric devices}
Real devices change their conductance by different amounts up and down, and the step depends on the state (soft bounds). Under
plain SGD this acts as an implicit penalty that pulls weights toward each device's symmetry point~\cite{gokmen2020,wu2025}.
\begin{itemize}
\item \textbf{Tiki-Taka (TTv1)}~\cite{gokmen2020} uses two arrays. The gradient is accumulated on an array $A$, and its columns
      are periodically transferred into the weight array $C$, with $W=\gamma A+C$. Asymmetry then acts as a useful decay on $A$
      instead of a bias on $W$.
\item \textbf{TTv2}~\cite{gokmen2021} inserts a digital buffer $H$ between the two arrays. $H$ low-pass filters the reads of $A$,
      and a pulse is sent to $C$ only when a threshold is crossed. TTv2 still needs every device's symmetry point to be
      subtracted by a reference.
\item \textbf{Chopped TTv2 (c-TTv2)} and \textbf{AGAD}~\cite{rasch2024,rasch2023arxiv} remove this sensitivity. A random sign
      modulates the inputs and the reads (c-TTv2). AGAD additionally replaces the static reference with a digital, on-the-fly
      reference taken from the recent history of $A$.
\item \textbf{Residual learning}~\cite{wu2025} gives convergence guarantees for this family.
\end{itemize}
IBM's open-source AIHWKIT~\cite{rasch2021aihwkit} implements these algorithms with fitted device presets, including the
Li-ECRAM and metal-oxide ECRAM models we use.

\subsection{Electrochemical three-terminal devices}
\begin{itemize}
\item \textbf{Lithium-ion synaptic transistors and ECRAM.} Fuller et al.~\cite{fuller2017} showed a lithium-ion synaptic
      transistor, followed by an ionic floating-gate array~\cite{fuller2019}. Tang et al.~\cite{tang2018} reported Li-WO$_3$ ECRAM
      with about a thousand symmetric, linear states and nanosecond pulses.
\item \textbf{Metal-oxide ECRAM.} Kim et al.~\cite{kim2019} made it CMOS-compatible.
\item \textbf{Protonic resistors.} Onen et al.~\cite{onen2022} reported nanosecond protonic programmable resistors with a
      phosphosilicate-glass electrolyte. In 2026, Xu et al.~\cite{xu2026} quantified cycle-to-cycle (3\%) and device-to-device
      (21\%) variability of BEOL-compatible protonic ECRAM.
\item \textbf{Open-loop programmable arrays.} Chen et al.~\cite{chen2023} demonstrated a $10\times10$ W/YSZ/WO$_x$ array with
      256 states.
\item \textbf{Organic devices.} Organic electrochemical synapses~\cite{vandeburgt2017} extend the family to low voltages.
\end{itemize}
In-situ training on ECRAM hardware has so far been shown on small arrays, for example a $3\times3$ array by Li et
al.~\cite{li2021ecram} and the $10\times10$ array of~\cite{chen2023}. Other memory technologies have gone further. Phase-change and
resistive arrays have trained or run networks at larger scale~\cite{prezioso2015,ambrogio2018,yao2020,wan2022}, and analog-AI
inference chips with on-chip periphery now report measured efficiencies of about 10~TOPS/W~\cite{ambrogio2023,legallo2023}.
